\documentclass{article}
\usepackage{amsmath, amssymb}
\usepackage{mathpartir}
\usepackage{geometry}
\geometry{margin=1in}

\title{Classifiers and Kinds: Definitions and Rules}
\author{}
\date{}

\begin{document}
\maketitle

\section{Classifier Structure}

Classifiers form a tree structure rooted at $\top$:

\begin{align*}
c &::= \top \mid c.n & \text{(classifiers)}
\end{align*}

where $c.n$ denotes the $n$-th child of classifier $c$.

\section{Kind Structure}

Kinds represent sets of classifiers with possible exclusions:

\begin{align*}
K &::= \emptyset \mid K \cup K \mid c \setminus [c_1, \ldots, c_n] & \text{(kinds)}
\end{align*}

The notation $c \setminus [\overline{c}]$ represents all classifiers in the subtree rooted at $c$, excluding the subtrees rooted at each $c_i$.

\paragraph{Abbreviations:}
\begin{itemize}
    \item $\mathsf{top} \triangleq \top \setminus []$ (all classifiers)
    \item $\mathsf{classifier}(c) \triangleq c \setminus []$ (subtree rooted at $c$)
\end{itemize}

\section{Classifier Subclassing}

The subclassing relation $c_1 \leq c_2$ holds when $c_1$ is in the subtree rooted at $c_2$:

\begin{mathpar}
\inferrule*[Right=SC-Refl]
{ }
{c \leq c}
\end{mathpar}

\begin{mathpar}
\inferrule*[Right=SC-Parent]
{c_1 \leq c_2}
{c_1.n \leq c_2}
\end{mathpar}

\section{Classifier Disjointness}

Two classifiers are disjoint ($c_1 \bot c_2$) when they are in different subtrees:

\begin{mathpar}
\inferrule*[Right=CD-Base]
{n \neq m}
{p.n \bot p.m}
\end{mathpar}

\begin{mathpar}
\inferrule*[Right=CD-Left]
{c_1 \bot c_2}
{c_1.n \bot c_2}
\end{mathpar}

\begin{mathpar}
\inferrule*[Right=CD-Right]
{c_1 \bot c_2}
{c_1 \bot c_2.m}
\end{mathpar}

\section{Kind Disjointness}

Two kinds are disjoint ($K_1 \bot K_2$) when they have no common classifiers:

\begin{mathpar}
\inferrule*[Right=KD-EmptyL]
{ }
{\emptyset \bot K}
\end{mathpar}

\begin{mathpar}
\inferrule*[Right=KD-EmptyR]
{ }
{K \bot \emptyset}
\end{mathpar}

\begin{mathpar}
\inferrule*[Right=KD-UnionL]
{K_1 \bot K \\ K_2 \bot K}
{K_1 \cup K_2 \bot K}
\end{mathpar}

\begin{mathpar}
\inferrule*[Right=KD-UnionR]
{K \bot K_1 \\ K \bot K_2}
{K \bot K_1 \cup K_2}
\end{mathpar}

\begin{mathpar}
\inferrule*[Right=KD-AbsurdL]
{r_1 \in \overline{e_1}}
{r_1 \setminus \overline{e_1} \bot r_2 \setminus \overline{e_2}}
\end{mathpar}

\begin{mathpar}
\inferrule*[Right=KD-AbsurdR]
{r_2 \in \overline{e_2}}
{r_1 \setminus \overline{e_1} \bot r_2 \setminus \overline{e_2}}
\end{mathpar}

\begin{mathpar}
\inferrule*[Right=KD-Root]
{r_1 \bot r_2}
{r_1 \setminus \overline{e_1} \bot r_2 \setminus \overline{e_2}}
\end{mathpar}

\begin{mathpar}
\inferrule*[Right=KD-ExclL]
{r_1 \in \overline{e_2}}
{r_1 \setminus \overline{e_1} \bot r_2 \setminus \overline{e_2}}
\end{mathpar}

\begin{mathpar}
\inferrule*[Right=KD-ExclR]
{r_2 \in \overline{e_1}}
{r_1 \setminus \overline{e_1} \bot r_2 \setminus \overline{e_2}}
\end{mathpar}

Here, $r \in \overline{e}$ means $\exists e_i \in \overline{e}.\ r \leq e_i$ (containment via superclass).

\section{Kind Intersection}

The intersection relation $K_1 \cap K_2 = R$ computes the common classifiers:

\begin{mathpar}
\inferrule*[Right=Int-EmptyL]
{ }
{\emptyset \cap K = \emptyset}
\end{mathpar}

\begin{mathpar}
\inferrule*[Right=Int-EmptyR]
{ }
{K \cap \emptyset = \emptyset}
\end{mathpar}

\begin{mathpar}
\inferrule*[Right=Int-UnionL]
{K_1 \cap K = R_1 \\ K_2 \cap K = R_2}
{(K_1 \cup K_2) \cap K = R_1 \cup R_2}
\end{mathpar}

\begin{mathpar}
\inferrule*[Right=Int-UnionR]
{K \cap K_1 = R_1 \\ K \cap K_2 = R_2}
{K \cap (K_1 \cup K_2) = R_1 \cup R_2}
\end{mathpar}

\begin{mathpar}
\inferrule*[Right=Int-SubL]
{r_1 \leq r_2}
{(r_1 \setminus \overline{e_1}) \cap (r_2 \setminus \overline{e_2}) = r_1 \setminus (\overline{e_1} \mathbin{+\mkern-8mu+} \overline{e_2})}
\end{mathpar}

\begin{mathpar}
\inferrule*[Right=Int-SubR]
{r_2 \leq r_1}
{(r_1 \setminus \overline{e_1}) \cap (r_2 \setminus \overline{e_2}) = r_2 \setminus (\overline{e_1} \mathbin{+\mkern-8mu+} \overline{e_2})}
\end{mathpar}

\begin{mathpar}
\inferrule*[Right=Int-Disj]
{r_1 \bot r_2}
{(r_1 \setminus \overline{e_1}) \cap (r_2 \setminus \overline{e_2}) = \emptyset}
\end{mathpar}

\section{Kind Subtraction}

The subtraction relation $K_1 - K_2 = R$ computes the classifiers in $K_1$ but not in $K_2$:

\begin{mathpar}
\inferrule*[Right=Sub-EmptyL]
{ }
{\emptyset - K = \emptyset}
\end{mathpar}

\begin{mathpar}
\inferrule*[Right=Sub-UnionL]
{K_1 - K = R_1 \\ K_2 - K = R_2}
{(K_1 \cup K_2) - K = R_1 \cup R_2}
\end{mathpar}

\begin{mathpar}
\inferrule*[Right=Sub-EmptyR]
{ }
{(r_1 \setminus \overline{e_1}) - \emptyset = r_1 \setminus \overline{e_1}}
\end{mathpar}

\begin{mathpar}
\inferrule*[Right=Sub-UnionR]
{(r_1 \setminus \overline{e_1}) - K_1 = R_1 \\ R_1 - K_2 = R_2}
{(r_1 \setminus \overline{e_1}) - (K_1 \cup K_2) = R_2}
\end{mathpar}

\begin{mathpar}
\inferrule*[Right=Sub-Tree]
{ }
{(r_1 \setminus \overline{e_1}) - (r_2 \setminus []) = r_1 \setminus (r_2 :: \overline{e_1})}
\end{mathpar}

\paragraph{Exclusion handling:} For $(r_1 \setminus \overline{e_1}) - (r_2 \setminus (a :: \overline{e_2}))$, using the identity $A - (B - C) = (A - B) \cup (A \cap B \cap C)$:

\begin{mathpar}
\inferrule*[Right=Sub-ExclAbsurd]
{a < r_2}
{(r_1 \setminus \overline{e_1}) - (r_2 \setminus (a :: \overline{e_2})) = r_1 \setminus \overline{e_1}}
\end{mathpar}

\begin{mathpar}
\inferrule*[Right=Sub-ExclIrrelR]
{a \bot r_2 \\ (r_1 \setminus \overline{e_1}) - (r_2 \setminus \overline{e_2}) = R}
{(r_1 \setminus \overline{e_1}) - (r_2 \setminus (a :: \overline{e_2})) = R}
\end{mathpar}

\begin{mathpar}
\inferrule*[Right=Sub-ExclSubR]
{a \leq r_2 \\ a \leq r_1 \\ (r_1 \setminus \overline{e_1}) - (r_2 \setminus \overline{e_2}) = R}
{(r_1 \setminus \overline{e_1}) - (r_2 \setminus (a :: \overline{e_2})) = R \cup (a \setminus \overline{e_1})}
\end{mathpar}

\begin{mathpar}
\inferrule*[Right=Sub-ExclSubL]
{a \leq r_2 \\ r_1 < a}
{(r_1 \setminus \overline{e_1}) - (r_2 \setminus (a :: \overline{e_2})) = r_1 \setminus \overline{e_1}}
\end{mathpar}

\begin{mathpar}
\inferrule*[Right=Sub-ExclIrrelL]
{a \leq r_2 \\ r_1 \bot a \\ (r_1 \setminus \overline{e_1}) - (r_2 \setminus \overline{e_2}) = R}
{(r_1 \setminus \overline{e_1}) - (r_2 \setminus (a :: \overline{e_2})) = R}
\end{mathpar}

Here, $a < r$ denotes strict subclassing ($a \leq r$ and $a \neq r$).

\section{Subkinding}

The subkinding relation $K_1 \subset K_2$ holds when every classifier in $K_1$ is also in $K_2$:

\begin{mathpar}
\inferrule*[Right=SK-Sub]
{K_1 - K_2 = R \\ \mathsf{isEmpty}(R)}
{K_1 \subset K_2}
\end{mathpar}

where $\mathsf{isEmpty}(K)$ is defined inductively:

\begin{mathpar}
\inferrule*[Right=IsEmpty-Empty]
{ }
{\mathsf{isEmpty}(\emptyset)}
\end{mathpar}

\begin{mathpar}
\inferrule*[Right=IsEmpty-Absurd]
{r \in \overline{e}}
{\mathsf{isEmpty}(r \setminus \overline{e})}
\end{mathpar}

\begin{mathpar}
\inferrule*[Right=IsEmpty-Union]
{\mathsf{isEmpty}(K_1) \\ \mathsf{isEmpty}(K_2)}
{\mathsf{isEmpty}(K_1 \cup K_2)}
\end{mathpar}

\end{document}
